The \wikilink{s45--b013}{Artin theorem} assumes a finite group of
automorphisms. The following different statement allows an infinite
group but assumes an algebraic normal extension, in the sense of
\wikilink{s39--b002}{the normality characterization}.

\begin{definition}\label{automorphism-fixed-field}
For a group $G$ of automorphisms of a field $K$, its \emph{fixed field}
is $K^G=\{a\in K:\sigma(a)=a\text{ for all }\sigma\in G\}$.
It is a field because the property is preserved by sums, products,
and inverses of nonzero elements.
\end{definition}
\begin{theorem}\label{normal-fixed-fields}
Let $K/k$ be algebraic and normal, let $G=\operatorname{Aut}(K/k)$,
and put $F=K^G$. Then $K/F$ is normal and separable, and $F/k$ is
purely inseparable. Neither $K/k$ nor $G$ is assumed finite.
\end{theorem}
\begin{proof}\label{normal-fixed-fields-proof}
For $\alpha\in K$, its $G$-orbit is finite because its images are
among the finitely many roots of its minimal polynomial over $k$.
The polynomial $P_\alpha(X)=\prod_{\beta\in G\alpha}(X-\beta)$ has
coefficients fixed by every element of $G$ and hence coefficients in
$F$. Its roots are distinct. The minimal polynomial over $F$ divides
$P_\alpha$, so $\alpha$ is separable over $F$. Apply
\wikilink{s40--ed-algebraic-separability--separable-algebraic-extension}{the
elementwise definition}. Normality over $F$ follows from
\wikilink{s39--b010}{normality above an intermediate field}.

Now take $a\in F$ and a root $b$ of its minimal polynomial over $k$
in an algebraic closure. The map $k(a)\to k(b)$ sending $a$ to $b$
extends to an embedding of $K$ into that closure by
\wikilink{s37--b012}{embedding extension}, using its proof with this
prescribed initial embedding. Normality makes this embedding an
automorphism of $K$ over $k$. Since $a\in F$, it fixes $a$, so $b=a$.
Thus $a$ has only one conjugate, and
\wikilink{s43--ed-algebraic-inseparability--unique-conjugate-criterion}{the
unique-conjugate criterion} proves that $F/k$ is purely inseparable.
\end{proof}
\begin{remark}\label{fixed-fields-scope}
In characteristic zero, the conclusion gives $F=k$. In positive
characteristic, a purely inseparable normal extension can have no
nonidentity base-field automorphisms, so its fixed field may be all
of $K$. No formula $[K:F]=|G|$ for arbitrary infinite $G$ has been
used or asserted, and no topology on $G$ is needed for this theorem.
\end{remark}
